\documentclass[12pt, a4paper]{article}
\usepackage[utf8]{inputenc}
\usepackage[T1]{fontenc}
\usepackage[ngerman]{babel}
\usepackage{amsmath, amssymb, amsthm}
\usepackage{geometry}
\usepackage{graphicx}
\usepackage{xcolor}
\usepackage{tcolorbox}
\usepackage{enumitem}
\usepackage{listings}
\usepackage{fancyhdr}
\usepackage{titlesec}
\usepackage{mdframed}
\usepackage{tikz}
\usepackage{pgfplots}
\pgfplotsset{compat=1.18}

\geometry{a4paper, margin=2.5cm}

\definecolor{theoryblue}{RGB}{30, 90, 160}
\definecolor{exerciseorange}{RGB}{200, 80, 0}
\definecolor{solutiongreen}{RGB}{0, 120, 60}
\definecolor{hintgray}{RGB}{100, 100, 100}
\definecolor{codebg}{RGB}{245, 245, 245}

% Theory box
\tcbuselibrary{skins, breakable}
\newtcolorbox{theorybox}[1]{
  colback=theoryblue!5,
  colframe=theoryblue,
  fonttitle=\bfseries\color{white},
  title=#1,
  breakable
}

% Exercise box
\newtcolorbox{exercisebox}[1]{
  colback=exerciseorange!5,
  colframe=exerciseorange,
  fonttitle=\bfseries\color{white},
  title=Aufgabe #1,
  breakable
}

% Solution space box
\newtcolorbox{solutionspace}[1][6cm]{
  colback=white,
  colframe=black!20,
  fonttitle=\itshape\color{hintgray},
  title=Rechenraum,
  height=#1,
  breakable
}

% Hint box
\newtcolorbox{hintbox}{
  colback=solutiongreen!5,
  colframe=solutiongreen,
  fonttitle=\bfseries\color{white},
  title=Hinweis,
  breakable
}

% Julia code style
\lstdefinelanguage{Julia}{
  keywords={function, end, return, if, else, elseif, for, while, in, begin, let, do, try, catch, finally, using, import, export, mutable, struct, abstract, type},
  keywordstyle=\color{theoryblue}\bfseries,
  comment=[l]{\#},
  commentstyle=\color{hintgray}\itshape,
  stringstyle=\color{solutiongreen},
  morestring=[b]",
  sensitive=true
}
\lstset{
  language=Julia,
  backgroundcolor=\color{codebg},
  basicstyle=\ttfamily\small,
  breaklines=true,
  frame=single,
  rulecolor=\color{black!30},
  numbers=left,
  numberstyle=\tiny\color{hintgray}
}

\pagestyle{fancy}
\fancyhf{}
\rhead{\textcolor{hintgray}{\small Rocket Science Workbook}}
\lhead{\textcolor{hintgray}{\small \leftmark}}
\cfoot{\thepage}

\titleformat{\section}{\large\bfseries\color{theoryblue}}{}{0em}{}[\titlerule]
\titleformat{\subsection}{\normalsize\bfseries\color{exerciseorange}}{}{0em}{}


\begin{document}

\begin{center}
  {\LARGE \textbf{Kapitel 07: Wissenschaftliche Visualisierung}}\\[0.5em]
  {\large \textcolor{hintgray}{2D- und 3D-Plots, Phasenraum, Interpretation}}\\[0.3em]
  {\small \textcolor{hintgray}{Rocket Science Workbook — simulation-2D-jl}}
\end{center}
\vspace{1em}
\hrule
\vspace{1.5em}

\section{Welche Plots braucht ein Ingenieur?}

\begin{theorybox}{Pflicht-Plots für eine Raketensimulation}
\begin{enumerate}
  \item \textbf{Trajektorie}: $x$ vs. $y$ — Wo fliegt die Rakete?
  \item \textbf{Zustandsgrößen}: $v(t)$, $h(t)$, $m(t)$ — Wie entwickelt sich der Zustand?
  \item \textbf{Kräfte}: $F_T(t)$, $F_D(t)$, $F_g(t)$ — Was treibt die Rakete an, was bremst sie?
  \item \textbf{Phasenraum}: $v_y$ vs. $h$ — Zeigt Energieübergänge und Orbitinsertion
  \item \textbf{3D-Perspektive}: Trajektorie als 3D-Plot (auch wenn nur 2D simuliert)
\end{enumerate}
\end{theorybox}

\section{Der Phasenraum}

\begin{theorybox}{Was ist der Phasenraum?}
Der \textbf{Phasenraum} einer ODE $\dot{\mathbf{s}} = f(\mathbf{s})$ ist der Raum aller möglichen Zustände. Ein Plot von $v_y$ vs. $h$ zeigt:
\begin{itemize}
  \item Wie sich Höhe und Geschwindigkeit gegenseitig beeinflussen
  \item Wo Energie umgewandelt wird (potentiell $\leftrightarrow$ kinetisch)
  \item Ob die Simulation sich einem Fixpunkt (Orbit) oder einem Grenzwert nähert
\end{itemize}
\end{theorybox}

\section{Makie.jl für 3D-Visualisierung}

\begin{theorybox}{Warum GLMakie?}
\begin{itemize}
  \item Interaktive 3D-Plots mit Rotation und Zoom
  \item Animationen mit \texttt{@animate} oder \texttt{Observable}
  \item Hochwertige Grafiken für Publikationen und Präsentationen
\end{itemize}
\texttt{Plots.jl} ist für statische 2D-Plots ideal; \texttt{GLMakie.jl} für 3D und Interaktivität.
\end{theorybox}

\section{Rechenaufgaben}

\begin{exercisebox}{1 — Phasenraum interpretieren}
Ein Plot von $v_y$ (y-Achse) vs. $h$ (x-Achse) zeigt zwei charakteristische Phasen:
\begin{itemize}
  \item Phase A: $h$ steigt, $v_y$ steigt
  \item Phase B: $h$ steigt weiter, aber $v_y$ sinkt
\end{itemize}
\textbf{a)} Was passiert physikalisch in Phase A? (Triebwerk aktiv? Welche Kraft dominiert?)\\[0.5em]
\textbf{b)} Was passiert in Phase B? Ist das Triebwerk noch an?\\[0.5em]
\textbf{c)} Wann ist $v_y = 0$ bei maximalem $h$ (Apogäum)? Wo liegt dieser Punkt im Phasenraum?
\end{exercisebox}

\begin{solutionspace}[6cm]
\end{solutionspace}

\begin{exercisebox}{2 — Energiebilanz}
Die spezifische mechanische Energie ist:
\[
  E(t) = \frac{1}{2}\left(v_x^2 + v_y^2\right) + g_0\,h
\]
(spezifisch = pro kg)\\[0.5em]
\textbf{a)} Berechne $E(t=0)$ für eine Rakete die an der Startrampe steht.\\[0.5em]
\textbf{b)} Für LEO gilt: $E_\text{LEO} \approx \frac{1}{2}v_\text{LEO}^2 + g\,h_\text{LEO}$ mit $v_\text{LEO}=7800\,\text{m/s}$, $h_\text{LEO}=400\,000\,\text{m}$.\\
Berechne $E_\text{LEO}$ und interpretiere das Ergebnis.\\[0.5em]
\textbf{c)} Das Triebwerk liefert die Leistung $P = T \cdot v$. Warum steigt die Effizienz mit der Geschwindigkeit? (Stichwort: Oberth-Effekt)
\end{exercisebox}

\begin{solutionspace}[7cm]
\end{solutionspace}

\begin{exercisebox}{3 — 3D-Erweiterungsgedanke}
Unsere Simulation ist 2D. Für eine echte 3D-Simulation bräuchten wir:
\[
  \mathbf{s} = [x, y, z, v_x, v_y, v_z, m]^\top \quad (7\text{-dimensional})
\]
\textbf{a)} Welche zusätzliche Kraft/Größe käme in 3D hinzu, die in 2D fehlt? (Stichwort: Coriolis, Azimut)\\[0.5em]
\textbf{b)} Was ist der Unterschied zwischen einem \textbf{polare Orbit} (Inklinationswinkel $\approx 90°$) und einem \textbf{äquatorialen Orbit} ($0°$)?\\[0.5em]
\textbf{c)} Warum starten die meisten Raketen möglichst nah am Äquator (z.B. Cape Canaveral)?\\
\textit{Hinweis: Erdrotation!}
\end{exercisebox}

\begin{solutionspace}[6cm]
\end{solutionspace}

\section{Verbindung zum Julia-Template}

\begin{theorybox}{Was du jetzt implementierst}
In \texttt{template.jl} (Kapitel 07) erstellst du ein vollständiges Dashboard:
\begin{enumerate}
  \item \textbf{6-Panel Dashboard}: Trajektorie, $h(t)$, $v(t)$, $m(t)$, Kräfte, Phasenraum
  \item \textbf{3D-Plot} mit \texttt{GLMakie}: Trajektorie als 3D-Kurve (z=0), rotierbar
  \item \textbf{Animation}: Rakete bewegt sich entlang der Trajektorie
  \item \textbf{Energie-Plot}: $E_\text{kin}(t)$, $E_\text{pot}(t)$, $E_\text{total}(t)$
\end{enumerate}
Das ist dein Showcase — der Plot den du jemandem zeigst.
\end{theorybox}

\end{document}
